% Source: https://artofproblemsolving.com/community/c3h2497136_equation
% Stable identity: aops-topic-2497136-q01
The equation says $m(2m-1)(2n-1)=2022=2\cdot3\cdot337$. Here $2n-1$ is an odd positive divisor of $2022$, so $2n-1\in\{1,3,337,1011\}$ and correspondingly
\[m(2m-1)\in\{2022,\ 674,\ 6,\ 2\}.\]
The function $m(2m-1)$ is increasing for positive integers $m$, taking the values $1,6,15,28,\dots$.
\begin{itemize}
\item $m(2m-1)=6$ gives $m=2$.
\item $m(2m-1)=2$ has no solution ($1<2<6$).
\item $2m^2-m-674=0$ has discriminant $1+8\cdot674=5393$, and $2m^2-m-2022=0$ has discriminant $1+8\cdot2022=16177$; neither is a perfect square ($73^2<5393<74^2$, $127^2<16177<128^2$), so there are no integer solutions.
\end{itemize}
So $m=2$ and $2n-1=337$, i.e. the only solution is $\boxed{(m,n)=(2,169)}$.
